\chapter{La vulnerabilità Format String}\label{chapter:fmt_vuln}
In questo capitolo è presentato in che modo la vulnerabilità \emph{format string} viene introdotta con un utilizzo errato della \texttt{printf} e quali azioni è possibile compiere quando essa è presente.


\section{Uso corretto e uso errato di printf}\label{sec:printf_right_wrong}
La funzione \texttt{printf} può essere utilizzata in diversi modi, alcuni corretti e sicuri, ed altri errati e vulnerabili.
Questo aspetto si basa principalmente su chi ha il controllo della \emph{format string}: il programmatore oppure l'utilizzatore del programma.

\subsection{Il pattern sicuro: format string come costante}\label{subsec:right_fmt}
Un uso corretto prevede che, all'interno della \emph{format string}, ogni specificatore di formato corrisponda a un argomento effettivamente passato alla funzione e che i rispettivi tipi di dato siano tra loro compatibili.
Per rispettare questa regola, la \emph{format string} deve essere scritta staticamente nel codice sorgente dal programmatore, in modo che sia sotto il suo esclusivo controllo (Codice~\ref{lst:right_fmt}).
\begin{lstlisting}[
    language=C,
    gobble=4,
    caption={Esempio di uso sicuro di \texttt{printf}},
    label={lst:right_fmt}
]
    #include <stdio.h>

    int main(void) {
        char buf[64];
        
        if (fgets(buf, sizeof(buf), stdin) == NULL) {
            return 1;
        }
        
        printf("%s", buf);

        return 0;
    }
\end{lstlisting}

Tuttavia, il fatto che sia scritta staticamente non evita che possano esserci \emph{undefined behavior} o un \emph{leak} di informazioni. Il programmatore ha la responsabilità di controllare il numero e il tipo degli specificatori rispetto agli argomenti.

\subsection{Il pattern vulnerabile: format string controllata dall'utente}\label{subsec:wrong_fmt}
Un uso errato e vulnerabile consiste nell'avere la \emph{format string} sotto il completo controllo dell'utente (Codice~\ref{lst:wrong_fmt}).
Di conseguenza, un attaccante può rompere facilmente la regola secondo cui ogni specificatore deve corrispondere ad un argomento effettivamente passato alla \texttt{printf}.
\begin{lstlisting}[
    language=C,
    gobble=4,
    caption={Esempio di uso vulnerabile di \texttt{printf}},
    label={lst:wrong_fmt}
]
    #include <stdio.h>

    int main(void) {
        char buf[64];

        if (fgets(buf, sizeof(buf), stdin) == NULL) {
            return 1;
        }

        printf(buf);

        return 0;
    }
\end{lstlisting}
Come detto nella Sezione~\ref{sec:callconv}, se uno specificatore non corrisponde a nessun argomento passato alla \texttt{printf}, esso va ad agire su dati residui presenti sui registri o sullo stack, in base ai valori correnti degli offset di \texttt{va\_list}.
Questo permette all'attaccante di causare \emph{leak} di informazioni o addirittura di eseguire scritture in memoria.


\section{Information disclosure: individuare e leggere un valore con \%x}\label{sec:info_disclosure}
Prendiamo il codice vulnerabile appena mostrato e aggiungiamo una variabile locale \texttt{local\_secret} che viene allocata nello \emph{stack frame} della funzione \texttt{main} (Codice~\ref{lst:leak_demo_c}).
Mostriamo, tramite lo sfruttamento della vulnerabilità \emph{format string}, come è possibile vederne il contenuto senza che questo venga passato esplicitamente a \texttt{printf} --- operazione altrimenti non fattibile.
\begin{lstlisting}[
    language=C,
    gobble=4,
    caption={Esempio di programma vulnerabile con variabile locale \texttt{local\_secret}},
    label={lst:leak_demo_c}
]
    #include <stdio.h>

    int main(void) {
        long local_secret = 0x1337;

        char buf[64];

        if (fgets(buf, sizeof(buf), stdin) == NULL) {
            return 1;
        }

        printf(buf);

        return 0;
    }
\end{lstlisting}

Per sapere dove la variabile \texttt{local\_secret} viene collocata esattamente, è necessario analizzare il disassemblato (Codice~\ref{lst:disas_leak_demo}), ovvero il codice assembly ricavato dalla traduzione del codice macchina dell'eseguibile, prodotto dalla compilazione di Codice~\ref{lst:leak_demo_c}. Da questo è possibile capire lo spazio che il \texttt{main} riserva alle sue variabili locali, dove esattamente le posiziona e anche la struttura del suo \emph{stack frame}.
\begin{lstlisting}[
    language={},
    numbers=none,
    gobble=4,
    caption={Estratto del disassemblato di \texttt{main}: posizione di \texttt{local\_secret}},
    label={lst:disas_leak_demo}
]

    ; ...

    push   rbp
    mov    rbp, rsp
    sub    rsp, 0x60
    
    ; ...
    
    mov    QWORD PTR [rbp-0x58], 0x1337
    
    ; ...
    
    call   printf@plt
    
    ; ...
\end{lstlisting}
Lo spazio per le variabili locali viene allocato da ``\texttt{sub rsp, 0x60}'', ovvero 96 byte.
Il registro \texttt{rsp} (\emph{stack pointer}) subito prima della \texttt{sub} ha lo stesso valore del registro \texttt{rbp} (\emph{base pointer}). Di conseguenza ``\texttt{mov QWORD PTR [rbp-0x58], 0x1337}'' indica che \texttt{local\_secret} si trova nel secondo \emph{qword} (blocco da 8 byte) a partire dalla cima dello stack (\texttt{rsp}), poiché il primo \emph{qword} si troverebbe a \texttt{rbp-0x60}.

Lo stack pointer non viene più modificato prima della chiamata a \texttt{printf} che quindi si ritrova con lo stack nello stato illustrato in Figura~\ref{fig:stack_leak_demo}.

\begin{figure}[ht]
    \centering
    \begin{tikzpicture}[
        box/.style={
            draw, thick,
            minimum width=4.8cm,
            minimum height=1cm,
            anchor=south west,
            font=\ttfamily\small
        },
        emptybox/.style={
            draw, thick,
            minimum width=4.8cm,
            minimum height=1cm,
            anchor=south west
        },
        ptr/.style={-{Stealth[length=3mm]}, thick},
        addr/.style={font=\ttfamily\footnotesize, anchor=west, inner sep=3pt}
    ]

    % Blocchi: base (indirizzi alti) in basso, cima (indirizzi bassi) in alto
    \node[box] (ret)    at (0,0) {saved ret addr};
    \node[box] (srbp)   at (0,1) {saved rbp};
    \node[emptybox, minimum height=2cm] (gapmid) at (0,2) {};
    \node at (2.4,3) {\color{gray}$\vdots$};
    \node[box] (secret) at (0,4) {local\_secret = 0x1337};
    \node[emptybox] (topgap) at (0,5) {};

    % Puntatori a sinistra: freccia allo spigolo alto-sinistra del blocco
    \node[font=\ttfamily\small, anchor=east] (rsplabel) at (-1.9,6) {rsp};
    \draw[ptr] (rsplabel.east) -- (0,6);
    \node[font=\ttfamily\small, anchor=east] (rbplabel) at (-1.9,2) {rbp};
    \draw[ptr] (rbplabel.east) -- (0,2);

    % Indirizzi a destra: allo spigolo alto-destra del blocco
    \node[addr] at (4.8,6) {rbp-0x60};
    \node[addr] at (4.8,5) {rbp-0x58};

    % Etichette indirizzi alti/bassi + freccia verticale
    \node[font=\small] (low)  at (8.1,6.3) {indirizzi più bassi};
    \node[font=\small] (high) at (8.1,-0.3) {indirizzi più alti};
    \draw[ptr] (low) -- (high);

    \end{tikzpicture}
    \caption{Stato dello \emph{stack frame} di \texttt{main} al momento della chiamata a \texttt{printf}, con la posizione di \texttt{local\_secret}}
    \label{fig:stack_leak_demo}
    
\end{figure}

La \emph{format string} occupa da sola il primo registro disponibile, \texttt{\%rdi}; restano quindi cinque registri \emph{general-purpose} (\texttt{\%rsi}, \texttt{\%rdx}, \texttt{\%rcx}, \texttt{\%r8}, \texttt{\%r9}) che la \emph{format string} non ha richiesto esplicitamente, ma che \texttt{va\_arg} considera comunque disponibili: sono gli ``argomenti fantasma'' introdotti nella Sezione~\ref{sec:callconv}. I primi cinque specificatori di formato leggono dunque, uno per uno, il contenuto residuo di questi cinque registri.

Esauriti i registri, il sesto specificatore legge il primo argomento dallo stack, puntato da \texttt{overflow\_arg\_area} a partire dall'\texttt{rsp} di \texttt{main} al momento della chiamata --- lo stesso \texttt{rsp} da cui abbiamo appena misurato la posizione di \texttt{local\_secret}. Il settimo specificatore legge quindi il secondo \emph{qword} a partire da \texttt{rsp}, che è esattamente dove si trova \texttt{local\_secret}.

Abbiamo due modi per arrivare a \texttt{local\_secret} (Codice~\ref{lst:leak_interaction}):
\begin{itemize}
    \item tramite una sequenza di specificatori di formato \texttt{\%lx};
    \item tramite lo specificatore di formato \texttt{\%7\$lx}.
\end{itemize}
\begin{lstlisting}[
    language={},
    numbers=none,
    gobble=4,
    caption={Invio della \emph{format string} per il \emph{leak} di \texttt{local\_secret}},
    label={lst:leak_interaction}
]
    %lx %lx %lx %lx %lx %lx %lx
    55790cd9e2a1 fbad2288 55790cd9e2bc 0 55790cd9e2a0 0 1337

    %7$lx
    1337
\end{lstlisting}

La sequenza di specificatori \texttt{\%lx} consuma, uno per uno, gli slot disponibili --- i cinque registri, poi lo stack --- fino a raggiungere \texttt{local\_secret} al settimo, stampato come \texttt{1337}.

Lo specificatore \texttt{\%7\$lx}, usato da solo, referenzia direttamente la settima posizione senza che le prime sei compaiano esplicitamente nella \emph{format string}: è esattamente il caso dei ``buchi'' descritto nella Sezione~\ref{sec:funcvar_format}, che la specifica richiede di evitare \cite{man3printf}. Il meccanismo che risolve la posizione richiesta si comporta però come se, per ciascuna delle sei posizioni non referenziate, fosse stato comunque letto un valore con \texttt{\%lx}. Non è un caso che i due metodi diano lo stesso risultato: entrambi consumano, nello stesso ordine, la stessa sequenza di slot \emph{general-purpose} (i cinque registri, poi lo stack) prima di raggiungere il settimo.\footnote{Se una delle posizioni saltate corrispondesse invece a un valore \emph{floating-point}, questa equivalenza implicita non varrebbe più: il conteggio dei registri \emph{general-purpose} avanzerebbe comunque, ma quello dei registri vettoriali no, disallineando i due contatori (verificato empiricamente).}

Il principio appena dimostrato non dipende dallo specifico specificatore \texttt{\%lx}: qualunque conversione che non dereferenzi l'argomento (come \texttt{\%x} e \texttt{\%p}), leggerebbe lo stesso registro o la stessa posizione di stack, cambiando solo il modo in cui il valore viene stampato. Lo stesso disallineamento tra specificatori dichiarati e argomenti realmente forniti che abbiamo sfruttato per leggere permette, come vedremo nella prossima sezione, anche di scrivere in memoria.


\section{Scrittura in memoria con \%n: il problema della dimensione}


\section{Scritture parziali: \%hhn e \%hn per la scrittura byte-a-byte}
